\documentclass[11pt,a4paper]{article}

\usepackage[utf8]{inputenc}
\usepackage[T1]{fontenc}
\usepackage[ngerman]{babel}
\usepackage{lmodern}
\usepackage{microtype}
\usepackage{geometry}
\geometry{margin=2.6cm}

\usepackage{amsmath,amssymb,amsfonts,amsthm}
\usepackage{graphicx}
\DeclareGraphicsExtensions{.pdf,.png,.jpg}
\graphicspath{{./}{fig/}{figs/}{images/}{img/}}
\newcommand{\includegraphicssafe}[2][]{\IfFileExists{#2}{\includegraphics[#1]{#2}}{\fbox{\texttt{Bilddatei #2 nicht gefunden}}}}

\usepackage{setspace}
\usepackage{enumitem}
\usepackage{booktabs}
\usepackage{xcolor}
\usepackage{hyperref}
\hypersetup{
	colorlinks=true,
	linkcolor=blue!50!black,
	urlcolor=blue!50!black,
	citecolor=blue!50!black,
	pdftitle={Hat die Zweiweg-Lichtgeschwindigkeit den Blick auf die globale Struktur der Physik verstellt?},
	pdfauthor={},
	pdfsubject={Philosophie der Physik, Relativitätstheorie, Messbegriff, Topologie},
	pdfkeywords={Lichtgeschwindigkeit, Zweiweg-Lichtgeschwindigkeit, Synchronisation, Relativität, Topologie, Sagnac-Effekt, Philosophie der Physik}
}

\newtheorem{satz}{Satz}
\newtheorem{definition}{Definition}
\newtheorem{lemma}{Lemma}
\newtheorem{korollar}{Korollar}
\theoremstyle{remark}
\newtheorem{bemerkung}{Bemerkung}
\newtheorem{beispiel}{Beispiel}

\title{\textbf{Hat die Zweiweg-Lichtgeschwindigkeit den Blick auf die globale Struktur der Physik verstellt?}\\
	\large Thomas Hoffbauer -Eine philosophische Kritik am heutigen physikalischen Weltbild}
\author{}
\date{}

\begin{document}
	\maketitle
	
	\begin{abstract}
		Die moderne Physik beschreibt die Lichtgeschwindigkeit mit außerordentlicher Präzision und erhebt sie zum Grundpfeiler ihres Weltbilds. Doch wie genau wird diese Konstanz eigentlich experimentell fixiert? In der Praxis fast immer über die Zweiweg-Lichtgeschwindigkeit, also über Hin- und Rückweg eines Signals. Der vorliegende Essay argumentiert, dass diese operative Festlegung zwar methodisch vernünftig und empirisch erfolgreich ist, zugleich aber eine philosophische Nebenwirkung haben könnte: Sie privilegiert lokale, symmetrisierte und synchronisationsstabile Beobachtungen und rückt dadurch globale, orientierungsabhängige und topologische Strukturen an den Rand der Aufmerksamkeit. Die These lautet nicht, dass die Relativitätstheorie falsch sei. Vielmehr könnte das heutige physikalische Weltbild die lokale Konsistenz der Natur schärfer sehen als ihre globale Form. Damit verbindet sich die Frage, ob Phänomene wie Rotation, Schleifen, Holonomie, Sagnac-Effekt oder andere orientierte Restgrößen nicht auf eine unterbelichtete Schicht physikalischer Wirklichkeit verweisen: auf eine Physik der globalen Präsenz statt bloß der lokalen Messbarkeit.
	\end{abstract}
	
	\section*{Einleitung}
	
	Kaum eine Konstante ist so ikonisch wie die Lichtgeschwindigkeit. Sie steht für die Grenzstruktur der Natur, für die Relativität von Raum und Zeit, für die Einheit der elektromagnetischen Theorie und für die tiefgreifende Umformung unseres Wirklichkeitsbegriffs im 20.\ Jahrhundert. In populären Darstellungen erscheint sie oft als geradezu sakrosankte Größe: exakt, universal, unerschütterlich.
	
	Doch eine philosophisch nüchterne Frage wird erstaunlich selten gestellt: \emph{Wie} kommt diese Konstanz eigentlich in die Welt der Messung? Nicht in die Welt der Gleichungen, sondern in die Welt der Uhren, Spiegel, Laufzeiten und Experimentalprotokolle. Und noch grundlegender: Welche Aspekte der Natur werden durch die Art der Messung hervorgehoben, welche treten zurück?
	
	Die leitende Vermutung dieses Essays ist einfach, aber folgenreich. Die moderne Physik hat die Lichtgeschwindigkeit vor allem als \emph{Zweiweg-Lichtgeschwindigkeit} operationalisiert. Gemessen wird also in aller Regel nicht die einfache Laufzeit von A nach B, sondern die Zeit für A nach B und wieder zurück nach A. Dies ist praktisch, robust und methodisch gut begründet. Doch gerade diese Symmetrisierung könnte den Blick auf eine andere Seite der Wirklichkeit verstellt haben: auf globale, orientierte, topologische oder schleifenartige Strukturen, die sich nicht in lokalen Rücklaufwerten erschöpfen.
	
	Diese Überlegung ist keine Einladung zu einem billigen Anti-Einstein-Reflex. Die Relativitätstheorie gehört zu den erfolgreichsten Theorien, die je formuliert wurden. Die hier vorgetragene Kritik ist subtiler. Sie richtet sich nicht primär gegen Formeln, sondern gegen eine implizite philosophische Selbstverständlichkeit des gegenwärtigen Weltbilds: gegen die Annahme, dass das methodisch Bestimmbare bereits mit dem ontologisch Relevanten zusammenfällt.
	
	\section{Die operative Vernunft der modernen Physik}
	
	Die moderne Physik ist nicht nur ein System von Theorien, sondern auch eine Kultur der Messung. Ihre Stärke liegt wesentlich darin, dass sie Begriffe an Operationen bindet. Was ist Länge? Was ist Zeit? Was ist Gleichzeitigkeit? Was ist Geschwindigkeit? Die Antwort lautet nicht mehr: das, was die Anschauung nahelegt, sondern: das, was sich reproduzierbar messen lässt.
	
	Gerade darin liegt ein gewaltiger Fortschritt. Die Geschichte der Physik ist auch die Geschichte der Befreiung von naiver Evidenz. Die Erde fühlt sich ruhig an und bewegt sich doch. Feste Körper wirken starr und bestehen doch aus quantenmechanischen Feldern. Zeit scheint gleichmäßig zu fließen und hängt doch vom Bewegungszustand und von Gravitationspotentialen ab.
	
	Aber jeder methodische Fortschritt hat seinen Preis. Wenn ein Begriff an eine Operation gebunden wird, dann gewinnt er an Klarheit und verliert zugleich einen Teil seiner Offenheit. Er wird präziser, aber enger. Das ist nicht falsch; es ist der normale Preis wissenschaftlicher Disziplin. Problematisch wird es erst, wenn man vergisst, dass diese Präzisierung immer auch eine Auswahl ist.
	
	Die Fixierung der Lichtgeschwindigkeit über zweiwegige Messverfahren ist dafür ein besonders schönes Beispiel. Sie löst ein praktisches Problem --- die Synchronisation entfernter Uhren --- mit eleganter Konsequenz. Aber gerade dadurch könnte sie eine tiefere Frage unsichtbar machen: Ob die Natur in ihrer globalen Gestalt mehr enthält als das, was auf lokalen Hin-und-Zurück-Pfaden stabil bleibt.
	
	\section{Warum die Physik die Zweiweg-Lichtgeschwindigkeit bevorzugt}
	
	Die Präferenz für die Zweiweg-Lichtgeschwindigkeit ist kein intellektueller Fehler, sondern eine methodische Notwendigkeit. Wer die einfache Laufzeit eines Lichtsignals von A nach B messen will, braucht an beiden Orten Uhren. Diese Uhren müssen synchronisiert sein. Doch genau dafür werden wieder Signale benötigt, typischerweise Lichtsignale. So entsteht ein Zirkel: Um die Einweg-Laufzeit zu messen, braucht man eine Synchronisation; um eine Synchronisation zu definieren, braucht man bereits ein Modell der Signalgeschwindigkeit.
	
	Die Zweiweg-Messung vermeidet dieses Problem. Ein Signal startet an einem Ort, läuft zu einem Spiegel und kehrt zum Ausgangspunkt zurück. Gemessen wird die gesamte Zeitdifferenz auf derselben Uhr. Damit wird keine Fernsynchronisation benötigt. Das Ergebnis ist robust, reproduzierbar und frei von einem Teil der konventionellen Zusatzannahmen, die bei Einweg-Messungen unvermeidlich wären.
	
	Genau an diesem Punkt beginnt jedoch die philosophische Frage. Wenn die fundamentalste Messung der Lichtgeschwindigkeit immer schon eine \emph{symmetrisierte Rücklaufmessung} ist, dann ist nicht ausgeschlossen, dass sie vor allem jene Aspekte der Natur hervorhebt, die unter Symmetrisierung invariant bleiben. Das ist ein großer Unterschied. Denn eine Methode, die besonders gut darin ist, lokale Konsistenz sichtbar zu machen, ist nicht automatisch ebenso gut darin, globale Differenzen ans Licht zu bringen.
	
	Mit anderen Worten: Vielleicht misst die Zweiweg-Lichtgeschwindigkeit die Natur genau dort am besten, wo sie sich selbst gleich bleibt. Aber was ist mit jenen Strukturen, die erst in Orientierung, Umlauf, Richtung, Schleifenbildung oder globaler Einbettung sichtbar werden?
	
	\section{Lokale Messbarkeit und globale Präsenz}
	
	Die moderne Physik ist in weiten Teilen eine Theorie der lokalen Struktur. Ihre Gleichungen sind lokal formuliert, ihre Größen sind punkt- oder feldweise definiert, ihre Invarianzen werden in infinitesimalen oder lokal linearen Näherungen ausgedrückt. Selbst dort, wo globale Aspekte relevant werden, geschieht der Zugang häufig über lokale Daten.
	
	Das ist außerordentlich mächtig. Aber nicht jede Wahrheit der Natur ist lokal.
	
	Ein Berg ist nicht dieselbe Art von Objekt wie ein Hangstück. Ein Umlauf ist nicht dieselbe Art von Phänomen wie ein gerader Wegabschnitt. Eine Schleife enthält Information, die kein noch so genauer Blick auf ein infinitesimales Linienelement vollständig liefert. Das ist keine Mystik, sondern eine banale topologische Einsicht. Es gibt Eigenschaften, die erst dann auftauchen, wenn ein Weg sich schließt, wenn ein Raum nicht trivial zusammenhängend ist, wenn Orientierung oder Einfassung physikalisch relevant werden.
	
	Die Geschichte der Physik kennt viele Beispiele dafür. Der magnetische Fluss durch eine Schleife ist nicht bloß die Summe lokaler Streckeninformationen. Holonomie ist keine Eigenschaft eines einzelnen Punktes, sondern eines Umlaufs. Der Aharonov-Bohm-Effekt erinnert daran, dass ein global eingeschlossener Zusammenhang physikalisch real sein kann, auch wenn entlang des Weges lokal nichts Spektakuläres geschieht. Der Sagnac-Effekt wiederum zeigt, dass gegenläufige Umläufe auf einer rotierenden Struktur nicht äquivalent sind, obwohl sie dieselbe geometrische Bahn benutzen.
	
	All diese Beispiele mahnen zur Vorsicht gegenüber einer Weltauffassung, die das Lokale für das Letzte hält. Vielleicht ist die Natur nicht nur eine Mannigfaltigkeit von Punkten mit Feldern, sondern ebenso eine Ordnung von Wegen, Umläufen, Rückbezügen und globalen Kopplungen.
	
	\section{Die verborgene Philosophie der Symmetrisierung}
	
	Die Zweiweg-Lichtgeschwindigkeit ist nicht nur eine technische Größe. Sie verkörpert eine bestimmte philosophische Haltung zur Wirklichkeit. Sie sagt implizit: Das physikalisch Relevante ist das, was auf geschlossenen Rücklaufpfaden stabil und reproduzierbar bleibt. Das ist eine vernünftige Haltung, aber eben eine Haltung.
	
	Symmetrisierung ist eine mächtige Methode. Sie glättet Störungen, mittelt Asymmetrien aus, erhöht die Robustheit und macht Invarianten sichtbar. Doch gerade deshalb kann sie jene Strukturen unterdrücken, die nur als Differenz, Orientierung oder Nichtäquivalenz existieren. Ein Hin-und-Zurück-Pfad ist kein neutraler Spiegel der Natur; er ist ein Filter.
	
	Dieser Gedanke ist leicht misszuverstehen. Gemeint ist nicht, dass die Physik sich geirrt habe, weil sie symmetrisiert. Gemeint ist vielmehr, dass jede Symmetrisierung erkenntnistheoretisch doppeldeutig ist. Sie offenbart das Stabile und verbirgt womöglich das Gerichtet-Gestaltete. Sie zeigt die Regel und verdeckt vielleicht den Zusammenhang, in dem die Regel gilt.
	
	Man könnte auch sagen: Die moderne Physik hat in bewundernswerter Weise gelernt, wie man sich gegen die Tücken lokaler Täuschung schützt. Vielleicht hat sie dabei aber zugleich verlernt, wie man globale Signaturen als eigenständige Realität ernst nimmt.
	
	\section{Einweg, Zweiweg und die Frage nach der Natur des Zusammenhangs}
	
	Aus rein experimenteller Sicht ist die Zweiweg-Lichtgeschwindigkeit der Goldstandard. Daraus folgt jedoch nicht notwendig, dass sie das ontologisch vollständigere Bild liefert. Ein Messprotokoll kann epistemisch privilegiert sein, ohne ontologisch exhaustiv zu sein.
	
	Hier liegt ein entscheidender Unterschied zwischen \emph{Messbarkeit} und \emph{Wirklichkeitsgehalt}. Die Physik ist zu Recht vorsichtig gegenüber allem, was nicht operational angebunden werden kann. Aber daraus entsteht schnell die Versuchung, die Grenze der Messmethode mit der Grenze der Sache selbst zu verwechseln.
	
	Die Einweg-Lichtgeschwindigkeit ist in der Standardtheorie eng mit Synchronisationskonventionen verbunden. Viele Physiker würden sagen: Genau deshalb besitzt nur die Zweiweg-Geschwindigkeit unmittelbaren physikalischen Gehalt. Das ist als methodische Aussage stark. Aber als philosophische Aussage ist es weniger selbstverständlich. Denn man könnte ebenso fragen, ob nicht gerade diese Konventionalität darauf hinweist, dass die lineare, offene, gerichtete Struktur von Signalfortpflanzung tiefer in globale Zusammenhänge eingebettet ist, die durch lokale Synchronisationsregeln nur partiell erfasst werden.
	
	Die eigentliche Frage wäre dann nicht: ``Wie schnell ist Licht wirklich?'' Sie wäre subtiler: \emph{Welche Struktur der Natur wird überhaupt sichtbar, wenn man Geschwindigkeit stets als lokal rückgekoppelte Größe definiert?}
	
	\section{Der Sagnac-Effekt als Stachel im lokalen Weltbild}
	
	Der Sagnac-Effekt nimmt in dieser Debatte eine besondere Stellung ein. Zwei Lichtsignale laufen auf einer rotierenden Plattform in entgegengesetzte Richtungen und treffen nicht gleichzeitig wieder ein. Das Resultat ist experimentell unstrittig und theoretisch gut verstanden. Doch philosophisch ist es bemerkenswert: Es macht deutlich, dass Orientierung, Umlauf und globale Einfassung realen physikalischen Unterschied erzeugen.
	
	Natürlich kann die etablierte Theorie dieses Phänomen behandeln. Es wäre billig, daraus unmittelbar eine Widerlegung der Relativitätstheorie konstruieren zu wollen. Aber gerade weil die Standardphysik den Effekt integrieren kann, ist er philosophisch so interessant. Er zeigt, dass die Natur an bestimmten Stellen nicht bloß lokal, sondern ausdrücklich \emph{umlaufbezogen} antwortet.
	
	In einem streng lokalen Weltbild wäre ein gegenläufiger Umlauf auf derselben Bahn zunächst etwas beinahe Gleiches. Doch die Rotation macht aus dieser scheinbaren Symmetrie eine physikalische Asymmetrie. Das Problem liegt also nicht auf der Ebene der lokalen Wegstücke, sondern in der globalen Einbettung des Weges.
	
	Damit wird eine prinzipielle Grenze sichtbar: Wer ausschließlich auf lokal symmetrisierte Größen blickt, kann leicht übersehen, dass die Natur auf geschlossene Wege anders reagiert als auf offene Strecken. Der Sagnac-Effekt ist in diesem Sinn weniger ein Angriff auf die moderne Physik als eine Erinnerung an etwas, das sie selbst weiß, aber in ihrem Selbstbild womöglich zu wenig betont: dass globale Orientierung keine Randerscheinung, sondern eine Grundkategorie physikalischer Beschreibung sein könnte.
	
	\section{Topologie als blinder Fleck einer metrisch dominierten Physik}
	
	Das 20.\ Jahrhundert war das Jahrhundert der Geometrisierung der Physik. Raum und Zeit wurden zur Raumzeit, Kräfte zu Krümmungen, Felder zu fundamentalen Größen, Symmetrien zu strukturbildenden Prinzipien. Das war ein enormer Fortschritt. Doch die Geometrie, die im Alltag des physikalischen Denkens dominiert, ist häufig metrisch und lokal differenzierbar gedacht.
	
	Topologie spielt zwar in vielen Spezialgebieten eine wachsende Rolle, von der Quantenmaterie bis zur Eichtheorie. Aber im allgemeinen philosophischen Selbstbild der Physik bleibt sie oft zweitrangig. Gemessen wird, vereinfacht gesagt, was sich lokal als Abstand, Zeitdauer, Energie oder Feldstärke ausdrücken lässt. Weniger präsent sind jene Aspekte, die sich durch geschlossene Wege, Einfassungen, Zusammenhangstypen, Orientierungsreste oder Holonomien ausdrücken.
	
	Gerade deshalb ist die Frage nach der Lichtgeschwindigkeit von solcher Tragweite. Wenn eine fundamentale Naturkonstante über einen rücklaufenden, geschlossenen Messmodus definiert wird, dann wäre es nicht erstaunlich, wenn der dabei bevorzugte Weltzugang metrisch-lokal bleibt, während topologisch-globale Eigenschaften nur in Spezialphänomenen auftauchen.
	
	Man könnte zugespitzt formulieren: Die moderne Physik versteht ausgezeichnet, \emph{wie schnell etwas lokal wieder bei sich ankommt}. Weniger klar ist, ob sie denselben begrifflichen Ernst dafür aufbringt, \emph{wie etwas global in der Welt anwesend ist}.
	
	\section{Von der Geschwindigkeit zur Präsenz}
	
	Vielleicht liegt das Problem tiefer als in der Frage nach einer bestimmten Messgröße. Vielleicht ist schon der Begriff der Geschwindigkeit selbst zu schmal, um die globale Natur von Signalfortpflanzung zu erfassen. Geschwindigkeit ist ein Quotient aus Weg und Zeit. Das ist sinnvoll für lineare Bewegungen, lokale Prozesse und kausale Abfolgen. Aber was, wenn die Welt nicht nur aus linearen Fortpflanzungen besteht, sondern auch aus Präsenzstrukturen?
	
	Mit \emph{Präsenz} ist hier nichts Esoterisches gemeint. Gemeint ist der globale Umstand, dass ein physikalischer Zusammenhang nicht nur lokal vorliegt, sondern als ganze Struktur wirksam ist. Ein Feld ist nicht bloß an einem Punkt; es ist verteilt. Eine Rotation ist nicht bloß ein lokaler Geschwindigkeitsvektor; sie ist ein Zustand des gesamten Umlaufsystems. Eine Schleife ist nicht bloß eine Summe kleiner Stücke; sie besitzt eine geschlossene Identität.
	
	Wenn man diese Perspektive ernst nimmt, könnte die Lichtgeschwindigkeit in ihrem üblichen physikalischen Sinn nur ein lokaler Grenzwert eines umfassenderen Zusammenhangs sein: einer globalen Kopplungsstruktur, einer topologischen Präsenzordnung oder einer Holonomie der Welt. Die Zweiweg-Lichtgeschwindigkeit würde dann nicht die ganze Wirklichkeit elektromagnetischer Signalfortpflanzung erfassen, sondern ihren lokal symmetrisierten Rücklaufwert.
	
	Das ist eine philosophische Möglichkeit, keine etablierte Theorie. Aber sie öffnet einen wichtigen Horizont. Sie erinnert daran, dass Naturbeschreibung nicht nur aus der Präzisierung von Zahlen besteht, sondern auch aus der Wahl dessen, was überhaupt als grundlegend gilt.
	
	\section{Kritik am heutigen Weltbild}
	
	Das gegenwärtige physikalische Weltbild ist außerordentlich erfolgreich, doch es leidet bisweilen an seiner eigenen Stärke. Weil es so viel erklären kann, neigt es dazu, seine bevorzugten Zugangsweisen mit der Struktur der Wirklichkeit selbst zu identifizieren. Daraus entstehen drei charakteristische Verkürzungen.
	
	\subsection*{Erste Verkürzung: Das Operative wird mit dem Ontologischen verwechselt}
	
	Was sich robust messen lässt, gilt schnell als das eigentlich Reale. Das ist verständlich, aber philosophisch zu kurz. Eine Messvorschrift selektiert; sie erschöpft nicht notwendig die Sache. Die Zweiweg-Lichtgeschwindigkeit ist operational hervorragend definiert. Daraus folgt nicht, dass die Welt selbst primär zweiwegig organisiert ist.
	
	\subsection*{Zweite Verkürzung: Das Lokale dominiert das Globale}
	
	Die Physik besitzt eine enorme Sensibilität für lokale Strukturen. Globale Differenzen erscheinen demgegenüber oft als Sonderfälle oder mathematische Verfeinerungen. Doch Natur könnte gerade dort wesentlich werden, wo Umläufe, Einfassungen und Nichttrivialitäten auftreten.
	
	\subsection*{Dritte Verkürzung: Symmetrie wird über Differenz gestellt}
	
	Die moderne Physik liebt Invarianten. Das ist einer ihrer größten Erfolge. Aber Invarianten sind nicht die ganze Geschichte. Eine Welt, die nur noch über das beschrieben wird, was unter Transformationen erhalten bleibt, verliert leicht die Aufmerksamkeit für das, was erst in orientierten Resten, Asymmetrien und gebrochenen Umläufen sichtbar wird.
	
	\section{Was eine neue philosophische Sensibilität leisten müsste}
	
	Eine ernsthafte Kritik am heutigen Weltbild darf nicht in die primitive Ablehnung seiner Erfolge umschlagen. Wer die Relativitätstheorie, die Quantenfeldtheorie oder die experimentelle Präzision moderner Messtechnik ignoriert, kritisiert nicht die Physik, sondern karikiert sie. Die Frage ist also nicht, ob wir das bestehende Gebäude einreißen sollten. Die Frage ist, ob wir seinen Grundriss zu eng interpretieren.
	
	Eine neue philosophische Sensibilität müsste daher drei Dinge leisten.
	
	Erstens müsste sie den Unterschied zwischen \emph{Messprotokoll} und \emph{Wirklichkeitsstruktur} offenhalten. Nicht, um sich von Messbarkeit zu verabschieden, sondern um zu vermeiden, dass methodische Tugenden zu ontologischen Dogmen werden.
	
	Zweitens müsste sie lokale und globale Beschreibungsebenen anders gewichten. Die Physik der Zukunft könnte dort weiterkommen, wo sie Schleifen, Holonomien, Rotationen, orientierte Differenzen und globale Kopplungen nicht als Randphänomene, sondern als konstitutive Größen ernst nimmt.
	
	Drittens müsste sie die Frage nach der \emph{topologischen Präsenz} rehabilitieren. Nicht alles, was wirkt, tut dies in Form eines lokalen Stoßes oder einer linearen Laufzeit. Vielleicht ist Natur an einigen entscheidenden Punkten eher durch globale Verbundenheit als durch lokale Übertragung gekennzeichnet.
	
	\section{Ein vorsichtiges Gegenbild}
	
	Nimmt man diese Überlegungen zusammen, dann entsteht kein Gegenentwurf zur modernen Physik im Sinne einer fertigen Alternativtheorie. Es entsteht eher ein Gegenbild im philosophischen Sinn.
	
	In diesem Gegenbild ist die Lichtgeschwindigkeit weiterhin eine fundamentale Größe. Aber sie erscheint nicht mehr als letzte und restlos transparente Wahrheit über Signalfortpflanzung, sondern als lokal symmetrisierte Manifestation eines tiefer liegenden Zusammenhangs. Offene Wege, geschlossene Umläufe und globale Orientierungen würden dann nicht bloß als Spezialfälle derselben Kinematik erscheinen, sondern als unterschiedliche ontologische Modi physikalischer Realität.
	
	Die leitende Frage wäre dann nicht mehr nur: ``Wie bewegen sich Signale?'' Sondern auch: ``Wie ist ein physikalischer Zusammenhang global anwesend?'' Aus dieser Perspektive könnten bekannte Phänomene --- Rotation, Sagnac-Effekt, Holonomie, topologische Quantisierung, globale Phasenbeziehungen --- in neuem Licht erscheinen.
	
	Vielleicht würde man am Ende feststellen, dass die bestehende Physik schon viel mehr davon enthält, als ihr philosophisches Selbstverständnis zugibt. Vielleicht aber auch, dass ein Teil ihres Weltbilds von einer stillen lokalen Voreingenommenheit geprägt ist.
	
	\section*{Schluss}
	
	Die Moderne hat uns gelehrt, dass die Wirklichkeit der Natur nicht identisch ist mit der Evidenz des Alltags. Vielleicht steht nun ein weiterer, kleinerer, aber nicht unbedeutender Schritt an: die Einsicht, dass auch die operative Evidenz der Physik nicht ohne Weiteres mit der ganzen Tiefe der Natur zusammenfällt.
	
	Die Zweiweg-Lichtgeschwindigkeit ist eine großartige Erfindung der Messvernunft. Sie hat es ermöglicht, eine stabile, konsistente und außerordentlich erfolgreiche Physik aufzubauen. Gerade deshalb sollte man sie philosophisch ernst nehmen --- ernst genug, um zu fragen, was sie sichtbar macht und was sie womöglich verdeckt.
	
	Vielleicht hat die Physik nicht zu wenig Präzision, sondern zu wenig Misstrauen gegenüber der Selektivität ihrer eigenen Präzision. Vielleicht sieht sie die lokale Wahrheit der Welt klarer als ihre globale Gestalt. Vielleicht ist die Frage nach der Lichtgeschwindigkeit deshalb noch nicht abgeschlossen, obwohl ihre Zahl mit unerreichter Genauigkeit feststeht.
	
	Nicht weil wir zu wenig messen. Sondern weil wir noch nicht ganz verstanden haben, was unsere Messungen von der Welt \emph{zeigen} --- und was sie uns übersehen lassen.
	
	\vspace{1.5em}
	\noindent\textbf{Kurz gesagt:} Die moderne Physik könnte die Konsistenz der Natur besser verstanden haben als ihre Präsenz.
	
\end{document}